\documentclass[11pt]{article}
\usepackage[margin=1in]{geometry}
\usepackage{amsmath,amssymb,amsthm,mathtools}
\usepackage{graphicx}
\usepackage[hidelinks]{hyperref}
\usepackage{microtype}

\newtheorem{theorem}{Theorem}
\newtheorem{lemma}{Lemma}
\newtheorem{remark}{Remark}
\newcommand{\A}{\mathcal A}
\newcommand{\Hh}{\mathcal H}
\newcommand{\M}{\mathcal M}
\newcommand{\Tr}{\operatorname{Tr}}
\newcommand{\rank}{\operatorname{rank}}
\newcommand{\N}{\mathbb N}

\title{No finite-rank operators do not force normality\\
\large A counterexample to Open Question 7.1 of arXiv:1009.3095}
\author{Automated mathematics run \texttt{fa\_banach\_001}, lane 2}
\date{9 August 2026}

\begin{document}
\maketitle

\noindent\textbf{Status.}
Candidate counterexample, likely valid; prepared for human review.

\section{Source question}

Lord and Sukochev define a spectral triple $(\A,H,D)$ as a countably
generated unital nondegenerate $*$-algebra of bounded operators on a
separable Hilbert space, together with a self-adjoint $D$ such that
$[D,a]$ is bounded for every $a\in\A$ and
$\langle D\rangle^{-1}=(1+D^2)^{-1/2}$ is compact.  If $n$ is the spectral
dimension and $\langle D\rangle^{-n}\in\M_{1,\infty}$, they consider
\[
 \Phi_\omega(S)=\Tr_\omega(S\langle D\rangle^{-n}),
 \qquad S\in\A''.
\]
Their Open Question 7.1 asks \cite[Section 7.2, p.~28]{LordSukochev}:
\begin{quote}
If $\A''$ contains no finite rank operators is $\Phi_\omega$ (the functional
in (1.1)) normal?
\end{quote}
As usual, ``contains no finite rank operators'' means that it contains no
\emph{nonzero} finite-rank operator.

\begin{figure}[ht]
  \centering
  \includegraphics[width=0.98\textwidth]{figures/open_problem_crop.png}
  \caption{Open Question 7.1 in the source paper, printed page 28.}
\end{figure}

\section{Idea of the proof}

The hypothesis only excludes projections whose Hilbert-space ranges are
finite-dimensional.  It does not exclude an atomic von Neumann algebra whose
minimal projections all have infinite multiplicity.  We represent
$\ell_\infty(\N)$ with precisely that multiplicity and choose $D$ with a
two-dimensional lattice spectrum.  The full operator
$\langle D\rangle^{-2}$ has logarithmically divergent trace and Dixmier trace
$1/2$, but the part over any fixed atom is trace class.  Hence the associated
functional vanishes on every atom and on every finite sum of atoms, while it
is nonzero at the identity.  Those finite sums increase strongly to the
identity, which is exactly a failure of normality.

\section{Construction and theorem}

Let $K=L=\ell_2(\N)$, with standard bases $(e_i)_{i\geq1}$ and
$(f_j)_{j\geq1}$, and put $H=K\otimes L$.  For $i\geq1$, let
\[
 P_i=|e_i\rangle\langle e_i|\otimes I_L.
\]
Define the unital $*$-algebra
\[
 \A=\left\{\lambda I_H+\sum_{i\in F}a_iP_i:
       \lambda,a_i\in\mathbb C,\ F\subset\N\text{ finite}\right\}.
\]
It is generated by $I_H,P_1,P_2,\ldots$.  Define the positive self-adjoint
operator $D$ by
\[
 D(e_i\otimes f_j)=(i+j)e_i\otimes f_j
\]
on its natural diagonal domain.

\begin{theorem}\label{thm:main}
The triple $(\A,H,D)$ is a spectral triple of spectral dimension $2$ with
$\langle D\rangle^{-2}\in\M_{1,\infty}$.  Its bicommutant contains no
nonzero finite-rank operator.  For every Dixmier trace $\Tr_\omega$ in the
conventions of the source paper, the positive tracial functional
\[
 \Phi_\omega(S)=\Tr_\omega(S\langle D\rangle^{-2}),
 \qquad S\in\A'',
\]
is not normal.  In fact,
\[
 \Phi_\omega(I_H)=\frac12,
 \qquad
 \Phi_\omega(P_1+\cdots+P_r)=0\quad(r\geq1).
\]
Consequently, the answer to Open Question 7.1 is negative.
\end{theorem}

\begin{proof}
Write $H_i=P_iH=\mathbb Ce_i\otimes L$.  Since every element of $\A$ is
scalar on each $H_i$, and since every $P_i$ belongs to $\A$, its commutant is
the algebra of bounded block-diagonal operators
\[
 \A'=\left\{\bigoplus_{i\geq1}X_i:X_i\in B(H_i),\
                         \sup_i\|X_i\|<\infty\right\}.
\]
The commutant of this algebra consists exactly of independently chosen
scalars on the blocks.  Therefore
\[
 \A''=\left\{\sum_{i\geq1}b_iP_i:(b_i)\in\ell_\infty(\N)\right\}
       \cong \ell_\infty(\N)\,\bar\otimes\,I_L.                 \tag{1}
\]
If $0\neq S=\sum_i b_iP_i$ belongs to $\A''$, choose $i$ with $b_i\neq0$.
The restriction of $S$ to the infinite-dimensional space $H_i$ is
$b_iI_{H_i}$, so $S$ has infinite rank.  Thus $\A''$ contains no nonzero
finite-rank operator.

The algebra $\A$ is countably generated, unital and nondegenerate.  Moreover
$D$ is diagonal in the same basis as every element of $\A$, so $[D,a]=0$ for
all $a\in\A$.  Each eigenvalue $m\geq2$ of $D$ has multiplicity $m-1$.
It follows that $\langle D\rangle^{-1}$ is compact and, for $p>0$,
\[
 \Tr(\langle D\rangle^{-p})
 =\sum_{m=2}^{\infty}(m-1)(1+m^2)^{-p/2}.
\]
The summand is asymptotic to $m^{1-p}$, so this series converges exactly when
$p>2$.  Thus the spectral dimension is $n=2$.

Set $T=\langle D\rangle^{-2}$.  Its eigenvalue
$(1+m^2)^{-1}$ has multiplicity $m-1$.  At the end of the $M$th complete
eigenspace shell there are
\[
 K_M=\sum_{m=2}^M(m-1)=\frac{M(M-1)}2
\]
eigenvalues, and their sum is
\[
 S_M=\sum_{m=2}^M\frac{m-1}{1+m^2}=\log M+O(1).                \tag{2}
\]
Since $\log(1+K_M)=2\log M+O(1)$, (2) gives
\[
 \frac{S_M}{\log(1+K_M)}\longrightarrow\frac12.              \tag{3}
\]
Between $K_M$ and $K_{M+1}$ the omitted part of one shell has total mass at
most $M/(1+(M+1)^2)=O(M^{-1})$, so (3) holds along all partial sums of the
decreasing eigenvalue sequence.  In particular $T\in\M_{1,\infty}$, is
Tauberian (Dixmier measurable), and the defining logarithmic mean has the
ordinary limit $1/2$.  Hence every normalized Dixmier trace under
consideration satisfies
\[
 \Tr_\omega(T)=\frac12.                                      \tag{4}
\]

For $r\geq1$, put $F_r=P_1+\cdots+P_r\in\A$.  The positive operator $F_rT$
is trace class because
\[
 \Tr(F_rT)=\sum_{i=1}^r\sum_{j=1}^{\infty}
                 \frac1{1+(i+j)^2}<\infty.                   \tag{5}
\]
Every Dixmier trace vanishes on trace-class operators, and therefore
$\Phi_\omega(F_r)=0$.  On the other hand, $F_r\uparrow I_H$ strongly, while
(4) yields $\Phi_\omega(I_H)=1/2$.  A positive normal functional preserves
such bounded increasing suprema.  Thus $\Phi_\omega$ is not normal.

Finally, (1) is abelian, so the positive functional $\Phi_\omega$ is
automatically tracial on $\A''$.  This completes the verification of every
hypothesis and the failure of the asserted conclusion.
\end{proof}

\section{Verification notes}

The counterexample does not exploit a noncommutative-extension ambiguity:
$\A''$ is abelian.  Nor does it rely on a special generalized limit: the
ordinary logarithmic eigenvalue mean exists, so (4) holds for every Dixmier
trace in the source classes.  Equation (5) is the decisive separation between
Hilbert-space rank and Dixmier mass: $F_r$ has infinite rank on $H$, as the
hypothesis permits, but it sees only finitely many summable rows of $T$.

A small numerical script in \texttt{code/check\_shell\_asymptotics.py}
checks the exact multiplicities, convergence of the ratios in (3), and
stabilization of finite-row sums.  It is a sanity check only; the proof is the
comparison argument above.

\section{Bounded novelty check and limitations}

The local run indexes were searched for arXiv:1009.3095, the paper title,
``no finite rank operators'', ``Dixmier trace normal'', and close variants;
no prior packet or attempt was found.  Exact-phrase and close-variant
web/arXiv searches were run for Open Question 7.1 and the normality
conjecture.  The OpenAlex record for the source (work W2101700424) listed 15
citing records through 9 August 2026, including duplicate arXiv/journal
versions.  Their titles were screened, and locally available full text was
queried for arXiv:1201.3661, 1210.3397, 2103.02067, and 2412.00628.  No later
paper explicitly answering Question 7.1, and no occurrence of this
construction, was found within those bounds.  This is a bounded novelty
audit, not a certification of priority.

The result answers only Open Question 7.1.  It does not characterize
normality and does not address the source paper's other questions about
measurability, generalized limits, or maximal symmetric subideals.  It uses
the deliberately broad spectral-triple definition printed in the source; a
stronger geometric axiom system might exclude the example and would require
a separately formulated question.

\section{Human-review recommendation}

Send to a specialist in noncommutative geometry or operator ideals.  The two
highest-value checks are: (i) the bicommutant computation (1), including the
infinite-multiplicity representation, and (ii) the normalization $1/2$ in
(3)--(4).  If those pass, the monotone-convergence obstruction is immediate.

\begin{thebibliography}{9}
\bibitem{LordSukochev}
S. Lord and F. Sukochev,
\emph{Measure Theory in Noncommutative Spaces},
SIGMA 6 (2010), 072, 36 pp.; arXiv:1009.3095.
\end{thebibliography}

\end{document}

